\section{Introduction}
\label{sec:intro}

Retrieval-augmented generation (RAG) increasingly uses learned
decisions inside inference. A model may decide whether to retrieve,
which evidence to rerank, when to reformulate a query, and whether to
escalate to a stronger generator
\citep{jeong2024adaptiverag,asai2024selfrag,han2026airrag,
zhai2025ragler,du2025iragkr,niu2026uarag}. These mechanisms improve the
attainable accuracy--cost frontier. Yet their gates and confidence
scores are learned from finite data and evaluated mainly by average
task performance. They do not bound how often the policy selected for
deployment will fail. They can also become unreliable when the query
mix or external knowledge changes.

We study this gap as a progressive inference problem. A RAG system has
an ordered set of computation rungs, from a cheap retrieval path to a
strong path with hybrid retrieval, reranking, verification, and a
larger generator. Each rung emits observable signals about retrieval
margin, support, citation coverage, freshness, and abstention. A learned
sufficiency model uses these signals to decide whether to stop or
acquire more evidence. The methodological question is how to retain
this adaptive computation while controlling the population risk of the
selected policy from finite calibration data.

\begin{figure*}[t]
\centering
\includegraphics[width=0.88\textwidth]{figures/fig_paradigm.png}
\caption{Three deployment paradigms for RAG knowledge services.
(a)~Point-estimate tuning deploys the cheapest configuration whose
\emph{measured} risk clears the target --- an uncorrected
multiple-testing procedure with no feasibility guarantee.
(b)~Distribution-free risk control bounds the risk of a \emph{fixed}
pipeline but neither searches a configuration space nor reduces cost,
and is silent under drift. (c)~\system\ combines search,
certification, and monitoring: an evidence contract is certified by a fixed-sequence
test over a policy space ordered by risk vouchers, and an anytime-valid
e-process monitor closes the loop after deployment.}
\label{fig:paradigm}
\end{figure*}

Current learning and selection practice does not answer this question.
Practitioners sample a development set, estimate each policy's quality,
and deploy the cheapest policy whose \emph{estimate} clears the target,
an approach we call point-estimate tuning
(Figure~\ref{fig:paradigm}a), shared
by workload-level pipeline optimizers \citep{russo2025abacus,
liu2025palimpzest,patel2024lotus} and by adaptive RAG routers and LLM
cascades \citep{jeong2024adaptiverag,chen2023frugalgpt,ong2024routellm}.
This selection is an uncorrected multiple-testing procedure. It favors
configurations whose measurement noise happens to be favorable, and the
problem worsens as the candidate set grows. As we quantify in
Section~\ref{sec:results}, state-of-practice optimizers deploy policies
whose true violation rate exceeds the declared target in up to 64\% of
calibration draws. Even a correctly tuned pipeline offers no protection after deployment,
when workload mix shifts, indexes fall behind the world, and knowledge
APIs degrade --- the knowledge drift problem that static evaluation
cannot see.

Distribution-free risk control addresses the statistical problem
(Figure~\ref{fig:paradigm}b): conformal
and Learn-then-Test methods \citep{angelopoulos2021ltt,
angelopoulos2024crc} can bound the violation rate of a \emph{fixed}
pipeline with finite-sample confidence, and recent work has applied
them to RAG abstention and answer sets \citep{li2024traq,kang2024crag}.
These methods treat the pipeline as given: they neither search a
configuration space nor reduce cost, and they are silent about what
happens when the calibrated distribution drifts. Conversely, the
systems that do search --- routers, cascades, Bayesian optimization ---
provide no feasibility statement. Search, certification, and
post-deployment monitoring therefore need to be designed together so
that cost optimization does not weaken the statistical guarantee.

We propose \system, which learns and certifies progressive RAG policies
(Figure~\ref{fig:paradigm}c). An \emph{evidence contract} states
admissible violation rates for answer quality, citation completeness,
evidence freshness, and latency, together with a confidence budget
$\delta$. Learned scores determine when to stop and which policies are
promising. A separate statistical layer determines whether the
selected policy satisfies the contract.

Three mechanisms support the method. First, empirical \emph{risk
vouchers} summarize the observed effect of retrieval truncation,
reranker removal, pre-filter pushdown, and generator downgrade. Their
bounds compose along rewrite chains without assumptions on rewrite
interactions. Vouchers order candidates but do not certify them.
Second, the rewrite lattice induces progressive policies. Per-rung
logistic models estimate whether the current answer is sufficient;
threshold policies, utility routers, group-conditional variants, and
human deferral determine the next action. A fixed-sequence
Hoeffding--Bentkus test certifies the least-cost policy in the accepted
prefix. Its validity is independent of the number of policies and of
the accuracy of the learned scores. Third, a restart-mixture e-process
monitors audited losses, detects drift with a delay bound independent
of the change time, localizes the change, and triggers recertification.
A geometric error schedule preserves validity across repeated epochs.

This work provides finite-sample guarantees for learned progressive RAG
over heterogeneous knowledge sources. Its main contributions are:
\begin{itemize}
\item We formulate progressive evidence acquisition as a learned
adaptive-computation problem under a multi-loss evidence contract
(Section~\ref{sec:prelim}).
\item We combine per-rung sufficiency learning, composable risk
vouchers, and a fixed-sequence selection algorithm. We prove
distribution-free contract validity for any score quality and give
conditions for cost consistency (Section~\ref{sec:method}).
\item We develop a restart-mixture drift detector whose delay bound
does not depend on the change time. It has constant-time updates,
remains valid with sampled audits, and provides a lifetime guarantee
across repeated recertification
(Section~\ref{sec:monitor}).
\item We evaluate the mechanism on four public benchmarks, four LLM
families, and two serving backends. The study includes score and policy
ablations, calibration scaling, dense candidate grids, audit rates,
strict and joint contracts, and controlled drift
(Sections~\ref{sec:setup}--\ref{sec:results}).
\end{itemize}

The remainder of this paper is organized as follows.
Section~\ref{sec:related} reviews related work.
Section~\ref{sec:prelim} formalizes evidence contracts and the
certified deployment problem. Section~\ref{sec:method} presents the
\system\ framework. Sections~\ref{sec:setup} and~\ref{sec:results}
report the experimental setup and results. Section~\ref{sec:discussion}
discusses implications, limitations, and threats to validity, and
Section~\ref{sec:conclusion} concludes.
